% !TeX root = ../main.tex
% ---------------------------------------------------------------
\section{Introduction}
\label{sec:introduction}
% ---------------------------------------------------------------

Lists of highly cited researchers carry weight well beyond bibliometrics. Universities cite them in promotion cases and press releases, ranking agencies use them as inputs, and funders and governments read them as signals of national standing. Since 2019 the most widely used of these lists has been the standardized citation metrics author database produced by Ioannidis and colleagues from Scopus, commonly known as the Stanford ``top 2\%'' list \citep{ioannidisStandardizedCitationMetrics2019a, ioannidisUpdatedSciencewideAuthor2020}. The database reports, for every author with at least five papers, six citation indicators and a composite $c$ that combines them \citep{ioannidisMultipleCitationIndicators2016}. Authors are listed if they are among the 100,000 highest by $c$ or in the top 2\% of their Science-Metrix subfield, with and without self-citations, for career-long citations and for citations received in a single recent year. The ninth version, released in 2026, lists about a quarter of a million researchers in each table.

The Stanford database answers several criticisms of earlier lists. It covers every field to a useful depth, separates self-citations, discounts co-authorship through the $h_m$ index \citep{schreiberCaseStudyArbitrariness2013}, and gives weight to single, first and last authorship. Its authors are explicit that it is not a measure of quality and that inclusion and exclusion should not be over-read. Three features nevertheless limit what it can show. First, $c$ is cumulative. Each of its six components grows with the number of papers and with career length, so the list favours prolific and senior researchers, and the single-year table only partly corrects this \citep{ioannidisStandardizedCitationMetrics2019a}. Second, every citation counts the same, whether it comes from a paper in a leading journal or from one in a venue that is rarely read, so the indicators are open to citation cartels and other forms of manipulation that the database's authors themselves have documented \citep{ioannidisGeneralizedViewSelfcitation2015, szomszorHowMuchToo2020}. Third, the database is built from a proprietary source, and its author profiles, field assignments and citation links cannot be inspected or reproduced outside Elsevier.

This paper describes an alternative built on open data and on a different idea of what an author-level indicator should measure. We start from the bipartite spectral ranking of sources and institutions of \citet{docampodomingoRecirculatingCardinalSpectral2026}, computed from OpenAlex \citep{priemOpenAlexFullyopenIndex2022}, which assigns every journal and institution in a field a size-independent influence $v$ with field mean one. Each work then inherits two scores: a publication score, from the influence of the source and institutions that produced it, and a citer score, the mean influence of the sources of the works that cite it. A researcher's score is the mean of these work scores over their papers, under full or fractional counting and with or without self-citations. The resulting indicators measure the typical standing of a researcher's papers rather than their volume, and they weight citations by where they come from.

We ask three questions. How many researchers can be scored at the volume threshold that the Stanford database uses, and how many of the Stanford-listed researchers can be found in OpenAlex? How closely do spectral scores agree with the Stanford indicators, and where they disagree, is the disagreement associated with field, country, self-citation or co-authorship? And how different would a top-2\% list be if it were drawn up on spectral scores rather than on $c$?

The paper proceeds as follows. Section~\ref{sec:spectral_ranking} sets out the spectral ranking and the work and researcher scores derived from it. Section~\ref{sec:data} describes the OpenAlex corpus, the volume threshold, the Stanford database and the matching of Stanford authors to OpenAlex. Section~\ref{sec:results} reports coverage, agreement and the composition of the two lists, and Section~\ref{sec:discussion} discusses what the comparison shows and what it does not.
